\section{Conclusion}\label{sec:conclusion}

LLM-written phishing turns a static detection problem into a streaming one. SLIDE treats it that way: a label-free sliced two-gate test on transformer representations decides when to buy labels, and a four-term robust objective decides how to use them. On a 28{,}000-message public stream that moves from human corpora to two LLM corpora and a paraphrase campaign, the test detected all 40 segment changes within one batch without a false alarm. With 2{,}000 labels SLIDE reached an MCC of 0.898, against 0.897--0.905 for the best detectors that consumed all 27{,}500 labels and 0.696--0.768 for detectors that never adapt. The value lies in the timing of the purchases: the same budget spent at random times lost 0.112 MCC, and oracle timing gained only 0.004. Random projections of the representation make that timing possible; the same test on the detector score alone lost 0.060 MCC. Of the four loss terms only the adversarial one improved accuracy measurably.

Two findings bound the result. Transformer heads, SLIDE included, confound LLM style with malicious intent and flag 6.8--27.0\% of LLM-written legitimate mail, and a 400-character benign prefix defeats every 128-token transformer detector we tested. Future work should pair the drift-adaptive transformer with a whole-message model and evaluate the policy on gradual and recurring drift in an operational stream; adaptive attackers who query the deployed detector belong in the next version of the threat model.
